% Generated draft: requires mathematical and source review.
\documentclass[11pt]{article}
\usepackage[margin=25mm]{geometry}
\usepackage{fontspec}
\setmainfont{Latin Modern Roman}
\usepackage{amsmath,amssymb,mathtools,unicode-math}
\setmathfont{Latin Modern Math}
\usepackage{enumitem,needspace,xurl,hyperref}
\setlength{\parindent}{0pt}
\setlength{\parskip}{5pt}
\setlength{\emergencystretch}{3em}
\newcommand{\N}{\mathbb{N}}
\newcommand{\Z}{\mathbb{Z}}
\newcommand{\Q}{\mathbb{Q}}
\newcommand{\R}{\mathbb{R}}
\newcommand{\C}{\mathbb{C}}
\newcommand{\F}{\mathbb{F}}
\newcommand{\A}{\mathbb{A}}
\newcommand{\PP}{\mathbb{P}}
\newcommand{\E}{\mathbb{E}}
\newcommand{\eps}{\varepsilon}
\newcommand{\dd}{\,\mathrm{d}}
\begin{document}
\section*{UnsolvedMath ID: 801034}
\textbf{Source problem code:} catalogue entry
\subsection{Definitions and mathematical statement}
For sets \(A,B\subseteq\N\), write \(A\leq_T B\) if \(A\) is computable with oracle \(B\), and write \(A\equiv_T B\) if both \(A\leq_T B\) and \(B\leq_T A\). The Turing degree of \(A\) is its equivalence class \(\deg_T(A)\). A set is computably enumerable (c.e.) if some algorithm enumerates exactly its elements. Let
\[
\mathcal D_{\mathrm{ce}}=\{\deg_T(A):A\subseteq\N\text{ is c.e.}\}.
\]
The supplied question is: ``Is every/any c.e. degree the degree of topological categoricity?''

The source record does not define ``degree of topological categoricity,'' specify the class of topological structures under consideration, or specify their effective presentations and the computational requirements on homeomorphisms. Consequently, a fully specified mathematical assertion cannot be recovered from this record alone.

To display the quantifier structure of the universal reading, let \(\mathcal C\) stand for the intended class of structures and let \(T(X,\mathbf d)\) stand for the intended assertion that \(\mathbf d\) is a degree of topological categoricity of \(X\). The question would then have the schema
\[
\forall\mathbf d\in\mathcal D_{\mathrm{ce}}\;\exists X\in\mathcal C\;T(X,\mathbf d).
\]
Here \(\mathcal C\) and \(T\) are unresolved placeholders, not definitions of a mathematical notion. The wording ``every/any'' also leaves the intended quantifier insufficiently explicit: if it asks whether there is any such c.e. degree, the corresponding schema is instead
\[
\exists\mathbf d\in\mathcal D_{\mathrm{ce}}\;\exists X\in\mathcal C\;T(X,\mathbf d).
\]
These readings are distinct and require source clarification.
\subsection{Short English statement}
The question asks whether c.e. Turing degrees can occur as degrees of topological categoricity. The supplied wording does not define topological categoricity or its effective setting, and ``every/any'' leaves the universal versus existential reading unresolved. A precise statement requires these specifications.
\subsection{Sources}
\begin{itemize}
\item Local UnsolvedMath registry, permanent ID \(801034\), source problem code \texttt{catalogue entry}, title \emph{Open Problems in Mathematical Logic --- Problem 46}. The supplied registry attributes the problem to Barmpalias, Gao et al., with proposed year 2026, and labels it \texttt{open}; this is supplied status metadata, not an independently verified mathematical assessment.
\item Barmpalias, Gao et al., \emph{Open Problems in Mathematical Logic} (2026), Problem 46. Supplied source URL: \url{https://arxiv.org/abs/2608.26628}. The provenance records access on 2026-10-16 and extraction method \texttt{arxiv\_latex}. These are supplied provenance details; the source has not been independently consulted or verified for this draft.
\end{itemize}
\end{document}
